%!TEX root = TheoretischePhysikC.tex

\section{Vorlesung: Spin}
Die Drehimpulsoperatoren erfüllen die folgende Vertauschungsrelationen:
\begin{equation}\label{eq:amcr}
	[\widehat{L}_1,\widehat{L}_2] = i\hbar \widehat{L}_3, \quad [\widehat{L}_2,\widehat{L}_3] = i\hbar \widehat{L}_1\quad\text{und}\quad [\widehat{L}_3,\widehat{L}_1] = i\hbar \widehat{L}_2.
\end{equation}
Es folgt, dass
\begin{equation}
	[\widehat{\mathbf{L}}^2, \widehat{L}_j]=0, \quad \text{für alle $j=1,2,3$,}
\end{equation}
wobei
\begin{equation}
	\widehat{\mathbf{L}}^2 \equiv \widehat{L}_1^2+\widehat{L}_2^2+\widehat{L}_3^2.
\end{equation}
Wir haben gesehen, dass die Bahndrehimpulsoperatoren $\widehat{\mathbf{L}} \equiv \widehat{\mathbf{r}}\times \widehat{\mathbf{p}}$ diese Kommutatoren erfüllen. Man sagt, dass $\widehat{\mathbf{r}}\times \widehat{\mathbf{p}}$ uns eine \emph{Darstellung} dieser Kommutatoren liefert. Eine wichtige Eigenschaft dieser Darstellung ist, dass den Gesamtdrehimpuls \emph{quantisiert} ist, im Sinne, dass die Eigenwerte des Gesamtdrehimpulsoperators $\widehat{\mathbf{L}}^2$ nur die Werte $\hbar^2 \ell(\ell+1)$ annehmen können, wobei $\ell = 0, \frac12, 1, \ldots$. Allerdings für die Bahndrehimpulsoperatoren $\widehat{\mathbf{r}}\times \widehat{\mathbf{p}}$ tauchen nur ganzzahlige Werte auf: $\ell = 0,1,2,\ldots$.

Eine faszinierende Frage bleibt offen: Gibt es Quantensysteme, die halb-ganzzahlige Eigenwerte aufweisen? Um diese Frage zu bejahen benötigen wir einen Hilbertraum $\mathcal{H}$ und drei Operatoren $\widehat{S}_j:\mathcal{H}\rightarrow \mathcal{H}$, $j=1,2,3$, die die Relationen (\ref{eq:amcr}) erfüllen:
\begin{equation}\label{eq:amcr}
	\eqbox{[\widehat{S}_1,\widehat{S}_2] = i\hbar \widehat{S}_3, \quad [\widehat{S}_2,\widehat{S}_3] = i\hbar \widehat{S}_1\quad\text{und}\quad [\widehat{S}_3,\widehat{S}_1] = i\hbar \widehat{S}_2.}
\end{equation}
Falls ein solcher Hilbertraum und solche Operatoren existieren, können wir zumindest feststellen, dass es keine \emph{mathematische} Hindernisse gibt die die Existenz von Quantensystemen mit halb-ganzzahligem gesamten Drehimpuls verhindern. Selbst wenn ein solches Quantensystem mathematisch existieren kann, ist es noch ein weiterer Schritt, festzustellen, ob ein solches System in der Natur vorkommen kann.

Es zeigt sich, dass Quantensysteme mit halbzahligen Eigenwerten für den gesamten Drehimpulsoperator tatsächlich existieren. Wir untersuchen nun den Fall $\ell = \frac12$. Der dazugehörige Hilbertraum ist durch $\mathcal{H} \equiv \mathbb{C}^2$ gegeben. Wir definieren dazu die Matrizen
\begin{equation}
	\widehat{S}_1 = \frac{\hbar}{2} \begin{pmatrix}
		0 & 1 \\ 1 & 0
	\end{pmatrix}, \quad \widehat{S}_2 = \frac{\hbar}{2} \begin{pmatrix}
		0 & -i \\ i & 0
	\end{pmatrix} \quad \text{und} \quad \widehat{S}_3 = \frac{\hbar}{2} \begin{pmatrix}
		1 & 0 \\ 0 & -1
	\end{pmatrix}.
\end{equation}
Der Gesamtdrehimpulsoperator $\widehat{\mathbf{S}}^2$ lautet
\begin{equation}
	\widehat{\mathbf{S}}^2 = \frac{3\hbar^2}{4} \mathbb{I}.
\end{equation}
Dieser Operator besitzt zwei (entartete) Eigenwerte, nämlich, $\frac{3\hbar^2}{4}$ und $\frac{3\hbar^2}{4}$. Vergleichen wir die Wirkung von $\widehat{\mathbf{S}}^2$ mit 
\begin{equation}
	\widehat{\mathbf{S}}^2|\ell,m\rangle = \hbar^2\ell(\ell+1)|\ell,m\rangle
\end{equation}
stellen wir fest, dass $\ell=\frac12$ sein muss. 

Bezüglich dieser Darstellung ist $\widehat{S}_3$ bereits diagonal. Das heisst, die Eigenbasis von $\widehat{S}_3$ darf mit $|\ell,m\rangle$ identifiziert werden:
\begin{equation}
	|\ell=\tfrac12,m=\tfrac12\rangle \equiv \begin{pmatrix}
		1 \\ 0
	\end{pmatrix}\quad\text{und}\quad 
	|\ell=\tfrac12,m=-\tfrac12\rangle \equiv\begin{pmatrix}
		0 \\ 1
	\end{pmatrix}.
\end{equation}

\section*{Zusammengesetzte Systeme}
Oft müssen wir bei der Modellierung eines quantenmechanischen Systems das System aus mehreren Teilsystemen aufbauen: aus einem Elektron und einem Proton wird dann ein Wasserstoff-Atom und so weiter; oder wir verfeinern unsere Beschreibung durch Hinzunahme weiterer Freiheitsgrade: Aus einem ``Elektron ohne Spin'', wird zusammen mit einer Beschreibung eines isolierten Spinfreiheitsgrads eine Theorie für Elektronen, die den Spin einschließt.

In all diesen Fällen verfahren wir nach der Regel, dass jeder Sorte quantenmechanischer
Systeme ein Hilbertraum zugeordnet wird. Sei also $\mathcal{H}_1$ der Hilbertraum der Systemsorte ``1'' und ebenso $\mathcal{H}_2$. Die Frage lautet also: Wie und in welchem Hilbertraum ``$\mathcal{H}_{\text{$1$\&$2$}}$'' sollen wir die Systemsorte ``1 \& 2'' von Systemen beschreiben, die aus einem Paar von Teilchen/Freiheitsgraden der Sorten 1 und 2 besteht? Die Antwort ist mit Hilfe durch das \emph{Tensorprodukt} $\mathcal{H}_1\otimes \mathcal{H}_2$ von $\mathcal{H}_1$ und $\mathcal{H}_2$ gegeben.

\begin{center}
\parbox[c][5cm][c]{14cm}{\begin{tcolorbox}[colback=banana,arc=0cm,frame hidden,enhanced]{Das Tensorprodukt $\mathcal{H}_1\otimes \mathcal{H}_2$ zweier Hilberträume $\mathcal{H}_1$ und $\mathcal{H}_2$ ist \emph{der} Hilbertraum, der Vektoren $\phi_1\otimes \phi_2$ mit $\phi_1\in \mathcal{H}_1$ und $\phi_2\in \mathcal{H}_2$ enthält, und von ihnen erzeugt wird, und in dem die Rechenregel 
\begin{equation}
	\langle \phi_1\otimes \phi_2|\psi_1\otimes\psi_2\rangle \equiv \langle\phi_1| \psi_1\rangle\langle\phi_2| \psi_2\rangle
\end{equation}
für Skalarprodukte gilt.
}\end{tcolorbox}}
\end{center}

Betrachten wir also als nächstes das Tensorprodukt von Operatoren. Für $A_1:\mathcal{H}_1 \rightarrow \mathcal{H}_1$ und $A_2:\mathcal{H}_2 \rightarrow \mathcal{H}_2$ definieren wir den Operator $A_1\otimes A_2:\mathcal{H}_1 \otimes\mathcal{H}_2 \rightarrow \mathcal{H}_1 \otimes\mathcal{H}_2$ durch die Gleichung
\begin{equation}
	\eqbox{(A_1\otimes A_2)(\phi_1\otimes \phi_2) \equiv (A_1\phi_1)\otimes (A_2\phi_2).}
\end{equation}
Da $A_1\otimes A_2$ ja ein linearer Operator sein soll, ist er damit auch auf allen
Linearkombinationen von Produktvektoren definiert.

Ein Tensorprodukt von Operatoren lässt sich natürlich auch durch seine Matrixelemente angeben. Direkt aus der Definition folgt
\begin{equation}
	\begin{split}
			[A_1\otimes A_2]_{jj',kk'} &\equiv \langle \phi_j\otimes \phi_{j'}|\left(A_1\otimes A_2\right) \phi_k\otimes\phi_{k'}\rangle \\
			&= \langle \phi_j\otimes \phi_{j'}| (A_1\phi_k)\otimes (A_2\phi_{k'})\rangle \\
			&= \langle \phi_j|A_1\phi_k\rangle \langle \phi_{j'}|A_2\phi_{k'}\rangle \\
			&= [A_1]_{j,k}[A_2]_{j',k'}.
	\end{split}
\end{equation}
Die Indizes an der Produktmatrix, sind also Paare von Indizes. Das Tensorprodukt einer $3\times 3$-Matrix mit einer $7\times 7$-Matrix ergibt also eine $21\times 21$-Matrix. Zur numerischen Verarbeitung solcher Produkte ist es nötig, aus den Paaren von Indizes wieder einfache Indizes zu machen, also alle möglichen Paare irgendwie durchzunummerieren.

Rechenregeln für das Tensorprodukt von Operatoren folgen meist leichter aus der Definition. Wir bekommen neben der Linearität in jedem Faktor, die für alle Tensorprodukte (auch von Vektoren) gilt:
\begin{align}
	(A_1\otimes A_2)(B_1\otimes B_2) &= A_1B_1\otimes A_2B_2 \\
	(A_1\otimes A_2)^\dag &= A_1^\dag \otimes A_2^\dag \\
	\text{$A_1\ge 0$ und $A_2\ge 0$} &\Rightarrow A_1\otimes A_2 \ge 0 \\
	\|A_1\otimes A_2\| &= \|A_1\|\| A_2\| \\
	\tr(A_1\otimes A_2) &= \tr(A_1)\tr(A_2).
\end{align}

Schließlich sollte noch erwähnt werden, dass alles, was hier über das Tensorprodukt aus zwei Faktoren gesagt wurde, sofort auf endlich viele Faktoren übertragbar ist.